\documentclass[11pt]{article}

\usepackage[a4paper, margin=1in]{geometry}

% Basic math packages
\usepackage{amsmath, amssymb, amsthm}
\usepackage{fontspec}      % XeLaTeX fonts

% Theorem environments
\newtheorem{theorem}{Theorem}
\newtheorem{lemma}{Lemma}
\newtheorem{proposition}{Proposition}
\newtheorem{corollary}{Corollary}

\theoremstyle{definition}
\newtheorem{definition}{Definition}

\theoremstyle{remark}
\newtheorem{remark}{Remark}

\begin{document}

\title{Practice Exercises}
\date{}
\maketitle

\begin{proof}
    we claim that
    
    \[
        P(n): (\sum_{k=0}^{n} a_k 10^k - \sum_{k=0}^{n} a_k) \equiv 0 \pmod 3
    \]

    or

    \[
        P(n): \sum_{k=0}^{n} a_k 10^k \equiv \sum_{k=0}^{n} a_k \pmod 3
    \]

    for the \textbf{base case} where $n_0 = 1$

    \[
    \begin{aligned}
    P(1) &= \sum_{k=0}^{1} a_k 10^k - \sum_{k=0}^{1} a_k  \\
    &= (a_0 + 10a_1) - (a_0 + a_1) \\
    &= 9a_1 \equiv 0 \pmod{3}.
    \end{aligned}
    \]

    then, we assume P(k) for arbitrary $(k \ge n_0)$, we prove $P(k+1)$

    \vspace{1em}

    every positive integer divisable by 3 will have all digits sum that's divisable by 3 as well, so

    \[P(k+1): \sum_{k=0}^{n} a_k 10^k + a_{k+1}10^{k+1} \equiv \sum_{k=0}^{n} a_k + a_{k+1} \pmod 3\]

    and since $P(k)$, our inductive hypothesis, on both side, already satisfied, we can ignore it, and focus on the term with coefficient of $k+1$

    \[a_{k+1}10^{k+1} \equiv a_{k+1} \pmod 3\]

    this is true, since we only focus on the sum of all digits, $a_{k+1}10^{k+1} = \{a_{k+1}\}\underbrace{00\ldots 0}_{k+1}$ --- which all it's digit sums up to $a_{k+1}$, same for $a_{k+1}$. since P(k+1) follows P(k), the \textbf{induction case} is satisfied, thus we can conclude the given relation is true \end{proof}

\end{document}
